\documentclass[11pt]{article}
\usepackage{mathblatt}
\begin{document}
\weit
\typeout{PROBE-BEGIN brueche-dezimalzahlen-e1-k1-s0-v2}
\begin{aufgabe}{\detokenize{brueche-dezimalzahlen-e1-k1-s0-v2}}
\begin{teile}
\teil Probe

\kreisdiagramm{/90,/90,/180}
\end{teile}
\end{aufgabe}
\typeout{PROBE-END brueche-dezimalzahlen-e1-k1-s0-v2}
\typeout{PROBE-BEGIN kreis-e3-k3-s4-v1}
\begin{aufgabe}{\detokenize{kreis-e3-k3-s4-v1}}
\begin{teile}
\teil Probe

\kreisdiagramm{Bus $162^\circ$/162, Rad $108^\circ$/108, Auto $90^\circ$/90}
\end{teile}
\end{aufgabe}
\typeout{PROBE-END kreis-e3-k3-s4-v1}
\typeout{PROBE-BEGIN wahrscheinlichkeit-e2-k2-s0-v1}
\begin{aufgabe}{\detokenize{wahrscheinlichkeit-e2-k2-s0-v1}}
\begin{teile}
\teil Probe

\kreisdiagramm{rot/1,blau/1,rot/1,gelb/1,rot/1}
\end{teile}
\end{aufgabe}
\typeout{PROBE-END wahrscheinlichkeit-e2-k2-s0-v1}
\typeout{PROBE-BEGIN brueche-dezimalzahlen-e1-k1-s0-v4}
\begin{aufgabe}{\detokenize{brueche-dezimalzahlen-e1-k1-s0-v4}}
\begin{teile}
\teil Probe

\kreisdiagramm{/90,/45,/45,/180}
\end{teile}
\end{aufgabe}
\typeout{PROBE-END brueche-dezimalzahlen-e1-k1-s0-v4}
\typeout{PROBE-BEGIN wahrscheinlichkeit-e2-k2-s0-v2}
\begin{aufgabe}{\detokenize{wahrscheinlichkeit-e2-k2-s0-v2}}
\begin{teile}
\teil Probe

\kreisdiagramm{grün/2,gelb/1,grün/1}
\end{teile}
\end{aufgabe}
\typeout{PROBE-END wahrscheinlichkeit-e2-k2-s0-v2}
\end{document}
